% =============================================================================
% Section 4: Empirics — Tables and Figures
% =============================================================================
\section{Empirical Evidence}\label{sec:empirics}

This section presents the key tables and figures that document the
defence--logistics duplex.  All values are computed from the
{\ISI} v1.0 data snapshot (EU-27, vintage~2024) and are identical
to those reported in Paper~2~\cite{isi_paper2}.

\subsection{Composite distribution}

\Cref{tab:key-summary} reports the summary statistics of the
composite distribution.  The {\ISI} ranges from 0.236 (France,
rank~27) to 0.517 (Malta, rank~1), with a population mean of
0.344 and standard deviation of 0.070.  The distribution is
compressed: the coefficient of variation is 0.203, the Gini
coefficient is 0.116, and 23 of 27~countries fall in the
moderately concentrated tier ($[0.25, 0.50)$).

% Table T01: Key summary statistics -- ISI composite, EU-27, 2024
\begin{table}[htbp]
\centering
\caption{Key summary statistics of the \ISI composite distribution, EU-27, vintage~2024, methodology~v1.0.}
\label{tab:key-summary}
\begin{tabular}{@{}lr@{}}
\toprule
{\textbf{Statistic}} & {\textbf{Value}} \\
\midrule
  Minimum (France, rank~27) & 0.23567126 \\
  Maximum (Malta, rank~1) & 0.51748148 \\
  Range & 0.28181022 \\
  Mean (population) & 0.34439144 \\
  Median & 0.34027984 \\
  Standard deviation (population) & 0.06992188 \\
  Coefficient of variation & 0.2030 \\
  IQR (P75 $-$ P25) & 0.09881618 \\
  P25 & 0.29179051 \\
  P75 & 0.39060669 \\
  Gini coefficient & 0.115463 \\
\addlinespace
  Highly concentrated ($\geq 0.50$) & 1 \\
  Moderately concentrated ($[0.25, 0.50)$) & 23 \\
  Mildly concentrated ($[0.15, 0.25)$) & 3 \\
\bottomrule
\end{tabular}
\end{table}

\subsection{Variance decomposition}

\Cref{tab:variance-decomp} reports the variance decomposition.
Defence contributes 35.10\% and logistics 34.26\% of total
composite variance.  The combined share is 69.37\%.

% Table T02: Variance decomposition of the ISI composite
\begin{table}[htbp]
\centering
\caption{Variance contribution decomposition of the \ISI composite.  Each axis's marginal contribution equals $(1/K^2)\sum_{k=1}^{K}\operatorname{Cov}(A_j, A_k)$, where $K=6$.  Own-variance and cross-covariance components are shown separately.}
\label{tab:variance-decomp}
\small
\begin{tabular}{@{}l rrrr@{}}
\toprule
\textbf{Axis} & \textbf{Own Var} & \textbf{Cross Cov} & \textbf{Marginal} & \textbf{\% Total} \\
\midrule
  Financial & 0.000039 & +0.000012 & 0.000050 & 1.03 \\
  Energy & 0.000064 & +0.000105 & 0.000169 & 3.46 \\
  Technology & 0.000249 & +0.000007 & 0.000255 & 5.23 \\
  Defence & 0.001414 & +0.000302 & 0.001716 & 35.10 \\
  Critical Inputs & 0.000412 & +0.000611 & 0.001023 & 20.92 \\
  Logistics & 0.000965 & +0.000710 & 0.001675 & 34.26 \\
\midrule
  Defence + Logistics & 0.002380 & +0.001012 & 0.003391 & 69.37 \\
  Remaining four axes & 0.000763 & +0.000735 & 0.001498 & 30.63 \\
\addlinespace
  Total & 0.003143 & +0.001746 & 0.004889 & 100.00 \\
\addlinespace
  Own-var share of total & \multicolumn{3}{r}{} & 64.28 \\
  Cross-cov share of total & \multicolumn{3}{r}{} & 35.72 \\
\bottomrule
\end{tabular}
\end{table}

\Cref{fig:variance-bar} provides a visual summary.  The dominance
of defence and logistics is immediate: the two bars on the right
dwarf the remaining four.

\begin{figure}[htbp]
  \centering
  % F01: Axis variance contributions (%) — horizontal bar chart
\begin{tikzpicture}
\begin{axis}[
  xbar,
  width=0.85\textwidth,
  height=6cm,
  bar width=10pt,
  xlabel={Contribution to composite variance (\%)},
  xlabel style={font=\small\sffamily},
  ytick=data,
  yticklabels={Financial, Energy, Technology, Critical Inputs, Logistics, Defence},
  yticklabel style={font=\small\sffamily},
  xticklabel style={font=\small\sffamily},
  xmin=0, xmax=40,
  enlarge y limits=0.15,
  nodes near coords,
  nodes near coords style={font=\tiny\sffamily, anchor=west},
  every node near coord/.append style={/pgf/number format/.cd,fixed,precision=1},
  axis lines*=left,
  grid=major,
  grid style={gray!20},
  major tick length=0pt,
]
\addplot[fill=ISIBlue!30, draw=ISIBlue] coordinates {
  (1.0326, 0)
  (3.4555, 1)
  (5.2258, 2)
  (20.9202, 3)
  (34.2646, 4)
  (35.1012, 5)
};
\end{axis}
\end{tikzpicture}

  \caption{Axis contributions to composite variance (\%).
    Defence (35.1\%) and logistics (34.3\%) together account for
    69.4\% of total composite variance.}
  \label{fig:variance-bar}
\end{figure}

\subsection{Axis dominance}

\Cref{tab:dominance-counts} reports the number of countries for
which each axis is the highest-scoring axis.

% Table T03: Axis dominance counts
\begin{table}[htbp]
\centering
\caption{Number of countries for which each axis is the highest-scoring axis.  Mean dominance gap: average difference between the highest and second-highest axis score among countries dominated by that axis.}
\label{tab:dominance-counts}
\begin{tabular}{@{}l r r l@{}}
\toprule
\textbf{Axis} & \textbf{Countries} & \textbf{Mean Gap} & \textbf{Share} \\
\midrule
  Financial & 0 & 0.00000000 & 0/27 (0.0\%) \\
  Energy & 0 & 0.00000000 & 0/27 (0.0\%) \\
  Technology & 0 & 0.00000000 & 0/27 (0.0\%) \\
  Defence & 16 & 0.18366014 & 16/27 (59.3\%) \\
  Critical Inputs & 0 & 0.00000000 & 0/27 (0.0\%) \\
  Logistics & 11 & 0.08847991 & 11/27 (40.7\%) \\
\midrule
  Defence + Logistics & 27 & {--} & 100.0\% \\
\bottomrule
\end{tabular}
\end{table}

\Cref{fig:dominance-share} visualises the binary split: defence
dominates 16~countries, logistics dominates~11, and no other axis
dominates any country.

\begin{figure}[htbp]
  \centering
  % F02: Axis dominance counts — bar chart
\begin{tikzpicture}
\begin{axis}[
  ybar,
  width=0.85\textwidth,
  height=6.5cm,
  bar width=18pt,
  ylabel={Number of countries},
  ylabel style={font=\small\sffamily},
  xtick=data,
  symbolic x coords={Financial, Energy, Technology, Defence,
                     {Crit.\ Inputs}, Logistics},
  xticklabel style={font=\small\sffamily, rotate=25, anchor=north east},
  yticklabel style={font=\small\sffamily},
  ymin=0, ymax=20,
  enlarge x limits=0.12,
  nodes near coords,
  nodes near coords style={font=\small\sffamily\bfseries, anchor=south},
  axis lines*=left,
  grid=major,
  grid style={gray!20},
  major tick length=0pt,
]
\addplot[fill=ISIBlue!30, draw=ISIBlue] coordinates {
  (Financial, 0)
  (Energy, 0)
  (Technology, 0)
  (Defence, 16)
  ({Crit.\ Inputs}, 0)
  (Logistics, 11)
};
\end{axis}
\end{tikzpicture}

  \caption{Number of countries dominated by each axis.
    Only defence and logistics appear; no country is dominated by
    financial, energy, technology, or critical inputs.}
  \label{fig:dominance-share}
\end{figure}

\subsection{Rank disruption under omission}

\Cref{tab:loo-disruption} reports the Spearman~$\rho$ and
Kendall~$\tau$ between the baseline ranking and each
leave-one-axis-out variant.

% Table T04: Leave-one-axis-out rank disruption
\begin{table}[htbp]
\centering
\caption{Leave-one-axis-out (LOO) rank disruption.  Spearman~$\rho$ and Kendall~$\tau$ measure rank-order agreement between the baseline and each LOO composite.  Lower values indicate greater rank disruption when the axis is removed.}
\label{tab:loo-disruption}
\begin{tabular}{@{}l rrrr@{}}
\toprule
\textbf{Excluded Axis} & $\boldsymbol{\rho}$ & $\boldsymbol{\tau}$ & \textbf{Mean} $|\Delta|$ & \textbf{Max} $|\Delta|$ \\
\midrule
  Financial & 0.992674 & 0.943020 & 0.67 & 2 \\
  Energy & 0.992063 & 0.943020 & 0.67 & 2 \\
  Technology & 0.968254 & 0.886040 & 1.26 & 7 \\
  Defence & 0.832723 & 0.675214 & 3.33 & 10 \\
  Critical Inputs & 0.941392 & 0.846154 & 1.78 & 8 \\
  Logistics & 0.926129 & 0.789174 & 2.07 & 7 \\
\bottomrule
\end{tabular}
\end{table}

\Cref{fig:loo-rho} plots $\rho$ for each omitted axis.  The
defence bar stands out: $\rho = 0.833$, substantially lower than
any other axis.

\begin{figure}[htbp]
  \centering
  % F03: LOO Spearman rho by excluded axis — bar chart
\begin{tikzpicture}
\begin{axis}[
  ybar,
  width=0.85\textwidth,
  height=6.5cm,
  bar width=18pt,
  ylabel={Spearman $\rho$ (baseline vs.\ LOO)},
  ylabel style={font=\small\sffamily},
  xtick=data,
  symbolic x coords={Financial, Energy, Technology, Defence,
                     {Crit.\ Inputs}, Logistics},
  xticklabel style={font=\small\sffamily, rotate=25, anchor=north east},
  yticklabel style={font=\small\sffamily},
  ymin=0.75, ymax=1.02,
  enlarge x limits=0.12,
  nodes near coords,
  nodes near coords style={font=\tiny\sffamily, anchor=south,
    /pgf/number format/.cd, fixed, precision=3},
  axis lines*=left,
  grid=major,
  grid style={gray!20},
  major tick length=0pt,
]
\addplot[fill=ISIBlue!30, draw=ISIBlue] coordinates {
  (Financial, 0.993)
  (Energy, 0.992)
  (Technology, 0.968)
  (Defence, 0.833)
  ({Crit.\ Inputs}, 0.941)
  (Logistics, 0.926)
};
\end{axis}
\end{tikzpicture}

  \caption{Spearman~$\rho$ between the baseline ranking and each
    leave-one-axis-out variant.  Removing the defence axis produces
    the largest rank disruption.}
  \label{fig:loo-rho}
\end{figure}
